% Nová příprava: 2026-10-28; 1. ročník SPŠ, etapa 13–24.
% Rozpis: 1_rocnik/tematicke_plany/Rozpis_vyuky_MAT_1_rocnik.tex, hodina 18.
% Původ: vlastní formulace výkladu a vlastní úlohy; metodická inspirace uvedenými zdroji.
% Kontrola, obsahové porovnání zdrojů a časy: planovani/DOKONCENI_PRIPRAV_13_24.md.
% Učebnice: ucebnice/01_Zakladni_poznatky/Mnoziny/Množiny.tex
% oddíly Rovnost množin, Podmnožiny, Množinové operace a Vennovy diagramy; dle tématu hodiny.
% Realisticky.cz: externi_zdroje/realisticky/01_Zakladni_poznatky/03_Mnoziny/010308_Intervaly/010308_Lekce.pdf
% Realisticky.cz: externi_zdroje/realisticky/01_Zakladni_poznatky/03_Mnoziny/010309_Mnoziny_shrnuti/010309_Lekce.pdf
\documentclass[11pt,a4paper]{article}
% Spolecny styl pro ucitelske pripravy MAT (Work/MAT)
% Verze 1.1, 2026-10-23
\NeedsTeXFormat{LaTeX2e}
\ProvidesPackage{priprava}[2026/10/23 v1.1 Sablona ucitelskych priprav MAT]

% Zakladni typografie a cestina
\RequirePackage[T1]{fontenc}
\RequirePackage[utf8]{inputenc}
\RequirePackage[czech]{babel}
\RequirePackage{lmodern}
\RequirePackage{microtype}
\RequirePackage{amsmath,amssymb}
\RequirePackage{xcolor}
\RequirePackage{enumitem}
\RequirePackage[most]{tcolorbox}
\RequirePackage{tabularx,booktabs,array}
\RequirePackage{fancyhdr}
\RequirePackage{needspace}
\RequirePackage{geometry}
\RequirePackage[hidelinks]{hyperref}

\geometry{a4paper,top=17mm,bottom=18mm,left=23mm,right=23mm}
\definecolor{notegray}{gray}{0.39}
\setlength{\parindent}{0pt}
\setlength{\parskip}{2.5pt plus 1pt minus 0.5pt}
\setlist[itemize]{topsep=2pt,itemsep=1pt,parsep=0pt,leftmargin=1.45em}
\setlist[enumerate]{topsep=2pt,itemsep=2pt,parsep=0pt,leftmargin=1.65em}
\setcounter{secnumdepth}{0}
\renewcommand{\arraystretch}{1.12}

% Zahlavi: vlevo rocnik | globalni tema, vpravo cislo: nazev hodiny
% Zapati zustava zcela prazdne.
\pagestyle{fancy}
\fancyhf{}
\renewcommand{\headrulewidth}{0.35pt}
\renewcommand{\footrulewidth}{0pt}
\setlength{\headheight}{14pt}
\newcommand{\nastavPripravu}[4]{%
  \fancyhead[L]{\small #1\enspace\textbar\enspace #2}%
  \fancyhead[R]{\small #3: #4}%
}

% Minimalisticke cernobile boxy. Kazda definice zvlast.
\tcbset{commonlesson/.style={enhanced,breakable=false,colback=white,colframe=black,
  colbacktitle=white,coltitle=black,boxrule=0.45pt,
  arc=0.8mm,outer arc=0.8mm,left=3mm,right=3mm,top=1.6mm,bottom=1.6mm,
  toptitle=1mm,bottomtitle=0.7mm,fonttitle=\bfseries,
  attach boxed title to top left={xshift=2.5mm,yshift=-1.15mm},
  boxed title style={colback=white,colframe=white,boxrule=0pt,left=1mm,right=1mm}}}
\newtcolorbox{zadaniBox}[1][]{commonlesson,title={Zadání},#1}
\newtcolorbox{prikladBox}[1][]{commonlesson,title={Příklad},#1}
\newtcolorbox{definiceBox}[1][]{commonlesson,title={Definice},#1}

% Spolecny box pro zadani a jeho reseni, bez mezery mezi castmi.
% Pouziti:
%   \begin{prikladReseniBox}
%     Zadani ulohy.
%     \tcblower
%     \begin{align*} ... \end{align*}
%   \end{prikladReseniBox}
% Spodni cast ma jemny sedy podklad a delici linku; lze v ni sazet
% vice radku matematiky. Obe casti maji jen jeden spolecny ramecek.
\tcbset{solutionlesson/.style={
  commonlesson,
  enhanced jigsaw,
  breakable,
  colbacklower=black!3,
  segmentation style={solid,draw=black!25,line width=0.35pt},
  fontlower=\small\color{notegray},
  before lower={\noindent\textbf{Řešení:}\enspace},
  leftlower=3mm,rightlower=3mm
}}
\newtcolorbox{prikladReseniBox}[1][]{solutionlesson,title={Příklad},#1}
\newtcolorbox{zadaniReseniBox}[1][]{solutionlesson,title={Zadání},#1}

% Pomucky pro autora pripravy: sedive poznamky; \reseni pro starsi pripravy
\newcommand{\poznamka}[1]{\par\begingroup\small\color{notegray}\itshape
  \textnormal{\textbf{Poznámka pro učitele.}}\enspace #1\par\endgroup}
\newcommand{\reseni}[1]{\par\begingroup\small\color{notegray}
  \textbf{Očekávané řešení:}\enspace #1\par\endgroup}
\newcommand{\blok}[1]{\needspace{5\baselineskip}\subsection*{#1}}
\newcommand{\vyrok}[1]{\uv{#1}}

\usepackage{tikz}
\usetikzlibrary{patterns,calc,angles,quotes}
\nastavPripravu{1. ročník}{Množiny}{18}{Operace s intervaly}
\begin{document}
\blok{Na ose platí stejné definice jako ve výčtu}
Intervaly už známe z hodiny 8. Dnes hledáme společné části a spojení
číselných množin. Základní množinou je $\mathbb R$, není-li uvedeno jinak.
Hranaté krajní body zapisujeme autorsky $\langle\,\rangle$, nikoli složenými závorkami.

Průnik vyžaduje obě podmínky současně, sjednocení alespoň jednu.
Nejprve nakresli obě množiny na jednu osu, potom urči výsledek a zvlášť
otestuj krajní body. Nekonečno není reálný krajní bod a nepíšeme u něj uzavřenou závorku.

\begin{prikladReseniBox}[title={P1: překrývající se intervaly}]
Urči $A\cap B$ a $A\cup B$, kde $A=\langle-3;2)$ a $B=(0;5\rangle$.
\tcblower
\begin{center}
\begin{tikzpicture}[x=.8cm,y=.6cm,color=black]
\draw[->] (-4,0)--(6,0) node[right] {$x$};
\foreach \x in {-3,0,2,5} {\draw (\x,.1)--(\x,-.1) node[below] {$\x$};}
\node[left] at (-4,1) {$A$}; \draw[very thick] (-3,1)--(2,1);
\fill (-3,1) circle(2pt); \draw[fill=white] (2,1) circle(2pt);
\node[left] at (-4,2) {$B$}; \draw[very thick] (0,2)--(5,2);
\draw[fill=white] (0,2) circle(2pt); \fill (5,2) circle(2pt);
\node[left] at (-4,3) {$A\cap B$}; \draw[very thick] (0,3)--(2,3);
\draw[fill=white] (0,3) circle(2pt); \draw[fill=white] (2,3) circle(2pt);
\end{tikzpicture}
\end{center}
Průnik je $(0;2)$: nula chybí v $B$, dvojka v $A$.
Sjednocení je $\langle-3;5\rangle$, protože od $-3$ do $5$ nechybí žádný bod
a oba vnější kraje jsou zahrnuté alespoň jedním intervalem.
\end{prikladReseniBox}

\begin{prikladReseniBox}[title={P2: dotyk v jediném bodě},breakable=false]
Porovnej dvojice $C=(-2;1\rangle$, $D=\langle1;4)$
a $E=(-2;1)$, $F=(1;4)$. Urči průniky a sjednocení.
\tcblower
$C\cap D=\{1\}$ a $C\cup D=(-2;4)$.
Jednoprvkovou množinu píšeme $\{1\}$, ne samotné číslo $1$.

$E\cap F=\varnothing$ a $E\cup F=(-2;1)\cup(1;4)$.
Bod $1$ chybí, proto sjednocení nesmíme změnit na celý $(-2;4)$.
Jediný chybějící bod už brání tomu, aby výsledek byl jedním intervalem.
\end{prikladReseniBox}

\begin{prikladReseniBox}[title={P3: interval není výčet celých čísel},breakable=false]
Urči průnik $I=\langle-2;3)$ s $\mathbb Z$ a s $\mathbb N$.
Lze $I\cap\mathbb Z$ označit jako interval $\langle-2;2\rangle$?
\tcblower
$I\cap\mathbb Z=\{-2,-1,0,1,2\}$ a $I\cap\mathbb N=\{1,2\}$.
První výsledek není reálný interval: například $\frac12$ je v $\langle-2;2\rangle$,
ale není v $\mathbb Z$. Složené a intervalové závorky mají jiný význam.
\end{prikladReseniBox}
\poznamka{Jádro P1–P3 a S1 a)–d): asi 38 min práce, organizace a závěr 7 min.
S1 e) je další nabídka podle tempa. U každého páru nech nejprve samostatný náčrt;
krajní bod nesmí rozhodovat jen podle „bližší“ závorky.}

\blok{Samostatné procvičení}
\begin{zadaniReseniBox}[title={S1: vždy urči průnik i sjednocení},breakable=false]
\begin{enumerate}[label=\alph*)]
\item $A=(-\infty;1\rangle$, $B=(-2;4)$.
\item $A=\langle0;2\rangle$, $B=(2;5\rangle$.
\item $A=(-3;-1)$, $B=\langle1;3\rangle$.
\item $A=(-5;5)$, $B=\langle-1;2\rangle$.
\item $A=(-\infty;0)$, $B=\langle0;+\infty)$.
\end{enumerate}
\tcblower
a) $(-2;1\rangle$; $(-\infty;4)$.
b) $\varnothing$; $\langle0;5\rangle$.
Dvojka nepatří do $B$, ale ve sjednocení je díky $A$.
c) $\varnothing$; $(-3;-1)\cup\langle1;3\rangle$.
d) $\langle-1;2\rangle$; $(-5;5)$, neboť $B\subseteq A$.
e) $\varnothing$; $\mathbb R$. Nula patří právě do druhé množiny.
\end{zadaniReseniBox}

\blok{Volitelná rezerva}
\begin{zadaniReseniBox}[title={R1: tři množiny a vzdálenost},breakable=false]
\begin{enumerate}[label=\alph*)]
\item Pro $A=(-\infty;3)$, $B=\langle-1;5\rangle$, $C=(1;+\infty)$
urči $(A\cap B)\cap C$ a $(A\cup B)\cap C$.
\item Čísla vzdálená od $2$ nejvýše o $3$ tvoří množinu $M$.
Zapiš ji intervalem a urči $M\cap(0;+\infty)$.
\end{enumerate}
\tcblower
a) $A\cap B=\langle-1;3)$, proto trojitý průnik $(1;3)$.
$A\cup B=(-\infty;5\rangle$, proto druhý výsledek $(1;5\rangle$.
b) $M=\langle-1;5\rangle$, neboť krajní body jsou $2-3=-1$ a $2+3=5$.
Hledaný průnik je $(0;5\rangle$. Použili jsme význam vzdálenosti z hodiny 9.
\end{zadaniReseniBox}
\poznamka{Rezerva 8–12 min; nevyžaduj obecný algoritmus nerovnic s absolutní hodnotou.}

\blok{Domácí procvičení}
\begin{zadaniReseniBox}[title={D1: 12–18 minut; kresli osy}]
\begin{enumerate}[label=\alph*)]
\item Pro $A=(-4;2\rangle$, $B=\langle-1;6)$ urči průnik a sjednocení.
\item Totéž pro $A=(-\infty;-1)$, $B=(-1;+\infty)$.
\item Totéž pro $A=\langle2;4\rangle$, $B=\langle4;7\rangle$.
\item Pro $I=(-2;3\rangle$ urči $I\cap\mathbb Z$ a doplněk $I$ v $\mathbb R$.
\item Pro $A=\langle-2;2\rangle$, $B=(0;5)$ a $C=\langle1;3\rangle$
urči $(A\cup B)\cap C$.
\end{enumerate}
\tcblower
a) $\langle-1;2\rangle$; $(-4;6)$.
b) $\varnothing$; $(-\infty;-1)\cup(-1;+\infty)$, tedy $\mathbb R\setminus\{-1\}$.
c) $\{4\}$; $\langle2;7\rangle$.
d) $\{-1,0,1,2,3\}$; $(-\infty;-2\rangle\cup(3;+\infty)$.
e) $A\cup B=\langle-2;5)$, průnik s $C$ je celé $C=\langle1;3\rangle$.
\end{zadaniReseniBox}
\end{document}
